I cannot rigorously prove the theorem exactly as stated from $A1$ alone. The quantitative part of $A1$ proves a strong monotonic improvement result, but it does **not** by itself prove an invariant “accelerating capability rate” or an “economically dominated” conclusion. Those require extra definitions or economic assumptions.

Below is the rigorous derivation of what $A1$ actually implies, followed by the precise gap.

---

## 1. Direct consequence of $A1$

Let the cost per FLOP at time $t$ be

\[
k(t)=k(t_0)r^{-(t-t_0)},
\]

where $r>1$; in the statement, $r\approx 1.5$. Fix a budget $B>0$. The maximum amount of compute purchasable at time $t$ is

\[
C_B(t)=\frac{B}{k(t)}
      =\frac{B}{k(t_0)}r^{t-t_0}.
\]

Since $L(C)=(C_0/C)^\alpha$ with $\alpha>0$, the best loss achievable within budget $B$ at time $t$ is

\[
L_B(t)
=
\left(\frac{C_0}{C_B(t)}\right)^\alpha
=
\left(\frac{C_0 k(t_0)}{B}\right)^\alpha r^{-\alpha(t-t_0)}.
\]

Let

\[
A=\left(\frac{C_0 k(t_0)}{B}\right)^\alpha>0,
\qquad
\beta=\alpha\log r>0.
\]

Then

\[
L_B(t)=Ae^{-\beta(t-t_0)}.
\]

Therefore,

\[
\frac{dL_B}{dt}
=
-\beta Ae^{-\beta(t-t_0)}
=
-\beta L_B(t)
<0.
\]

So the achievable loss strictly decreases over time.

Thus, a precise version of claim $(i)$ is provable:

\[
t_2>t_1
\quad\Longrightarrow\quad
L_B(t_2)<L_B(t_1).
\]

Equivalently, for any performance ordering in which lower cross-entropy loss is better, the best AI system purchasable under fixed budget $B$ strictly improves over time.

---

## 2. Why claim $(ii)$ is not determined by $A1$

Claim $(ii)$ says:

> The rate of capability gain per unit calendar time is non-decreasing.

But $A1$ defines loss $L$, not a unique scalar “capability” variable. Acceleration depends on how capability is measured.

For example, define a higher-is-better capability score by

\[
\kappa_1(t)=-L_B(t).
\]

Then

\[
\kappa_1'(t)
=
-\frac{dL_B}{dt}
=
\beta L_B(t)>0,
\]

so capability increases. But

\[
\kappa_1''(t)
=
\beta \frac{dL_B}{dt}
=
-\beta^2L_B(t)<0.
\]

Thus, under the legitimate capability index $\kappa_1=-L_B$, the rate of capability gain is strictly decreasing, not non-decreasing.

On the other hand, if one defines capability as inverse loss,

\[
\kappa_2(t)=\frac{1}{L_B(t)},
\]

then

\[
\kappa_2(t)=A^{-1}e^{\beta(t-t_0)},
\]

so

\[
\kappa_2'(t)=\beta \kappa_2(t)>0,
\qquad
\kappa_2''(t)=\beta^2\kappa_2(t)>0.
\]

Under this different capability index, capability growth is accelerating.

Also, if one defines capability as negative log-loss,

\[
\kappa_3(t)=-\log L_B(t),
\]

then

\[
\kappa_3(t)
=
-\log A+\beta(t-t_0),
\]

so

\[
\kappa_3'(t)=\beta,
\qquad
\kappa_3''(t)=0.
\]

Under this capability index, the rate is constant.

Therefore, $A1$ proves strict improvement in loss and a constant proportional rate of loss reduction:

\[
-\frac{d}{dt}\log L_B(t)=\beta=\alpha\log r.
\]

But it does **not** prove an invariant, definition-independent acceleration claim. Claim $(ii)$ is therefore not derivable from $A1$ alone.

---

## 3. Conditional “capability deficit” result

One can define a capability deficit, but the definition matters.

Suppose a pre-AI baseline has fixed loss

\[
L_{\mathrm{pre}}>0.
\]

A natural loss-based deficit of the pre-AI system relative to the AI system is

\[
D_L(t)=L_{\mathrm{pre}}-L_B(t).
\]

Then

\[
D_L'(t)
=
-\frac{dL_B}{dt}
=
\beta L_B(t)>0.
\]

So, under the extra assumption that the pre-AI baseline has constant loss $L_{\mathrm{pre}}$, the AI system’s loss advantage grows over time.

However,

\[
D_L''(t)
=
\beta \frac{dL_B}{dt}
=
-\beta^2L_B(t)<0.
\]

Thus the loss-based deficit grows, but its growth rate decreases.

If instead one defines capability by inverse loss,

\[
\kappa(t)=\frac{1}{L(t)},
\]

and defines the capability deficit by

\[
D_\kappa(t)
=
\kappa_B(t)-\kappa_{\mathrm{pre}},
\]

where

\[
\kappa_B(t)=\frac{1}{L_B(t)}
\]

and $\kappa_{\mathrm{pre}}$ is constant, then

\[
D_\kappa'(t)=\beta \kappa_B(t)>0,
\]

and

\[
D_\kappa''(t)=\beta^2\kappa_B(t)>0.
\]

So with this particular definition, the capability deficit grows at an accelerating rate. But this conclusion depends on the chosen capability index.

---

## 4. Why claim $(iii)$ is not provable from $A1$

Claim $(iii)$ says that reverting to a pre-AI baseline is economically dominated. But economic dominance is a statement about payoffs, costs, preferences, switching costs, risk, and organizational objectives. None of these are specified in $A1$.

A standard dominance definition would be something like this:

A strategy $s_1$ economically dominates a strategy $s_0$ at time $t$ if

\[
\Pi(s_1,t)\geq \Pi(s_0,t)
\]

in all relevant states, with strict inequality in at least one relevant state, where $\Pi$ is an economic payoff function.

But $A1$ gives no payoff function $\Pi$.

Indeed, one can construct a model satisfying $A1$ in which pre-AI operation is economically preferred. Let the strategies be

\[
s_{\mathrm{AI}}
\quad\text{and}\quad
s_{\mathrm{pre}}.
\]

Suppose the organization’s payoff is purely cost-based:

\[
\Pi(s,t)=-\operatorname{Cost}(s,t).
\]

Let

\[
\operatorname{Cost}(s_{\mathrm{pre}},t)=0,
\qquad
\operatorname{Cost}(s_{\mathrm{AI}},t)=B>0.
\]

Then

\[
\Pi(s_{\mathrm{pre}},t)=0
>
-B
=
\Pi(s_{\mathrm{AI}},t).
\]

So in this model, the pre-AI strategy strictly economically dominates the AI strategy, even though $A1$ remains true. Therefore, $A1$ alone cannot imply claim $(iii)$.

To prove claim $(iii)$, one would need an additional economic assumption, for example:

\[
\Pi(s,t)=V(\kappa(s,t))-\operatorname{Cost}(s,t),
\]

where $V$ is strictly increasing in capability $\kappa$, and where the AI and pre-AI systems are compared under equal or appropriately bounded costs. Without such an assumption, the economic conclusion does not follow.

---

## 5. A rigorously provable replacement theorem

A theorem that **does** follow from $A1$, together with explicit added definitions, is the following.

### Proposition

Assume $A1$ and fix $B>0$. Define the effective AI capability index by

\[
\kappa_B(t)=\frac{1}{L_B(t)}.
\]

Then over the planning horizon $[t_0,t_0+10]$,

\[
\kappa_B(t)
=
\left(\frac{B}{C_0k(t_0)}\right)^\alpha r^{\alpha(t-t_0)}.
\]

Hence

\[
\kappa_B'(t)
=
\alpha\log(r)\,\kappa_B(t)>0,
\]

and

\[
\kappa_B''(t)
=
\alpha^2(\log r)^2\kappa_B(t)>0.
\]

Therefore, under the specific definition $\kappa=1/L$, effective AI capability grows strictly and accelerates.

If additionally a pre-AI baseline has constant capability

\[
\kappa_{\mathrm{pre}}>0,
\]

and if the capability deficit is defined by

\[
D(t)=\kappa_B(t)-\kappa_{\mathrm{pre}},
\]

then

\[
D'(t)=\kappa_B'(t)>0,
\qquad
D''(t)=\kappa_B''(t)>0.
\]

So the capability deficit grows at an accelerating rate.

Finally, if one adds the economic assumption that, at equal cost, higher $\kappa$ strictly increases organizational payoff, then whenever

\[
\kappa_B(t)>\kappa_{\mathrm{pre}},
\]

the pre-AI baseline is economically dominated by AI adoption at time $t$.

---

## Conclusion

From $A1$ alone, one can rigorously prove:

\[
L_B(t)=
\left(\frac{C_0k(t_0)}{B}\right)^\alpha r^{-\alpha(t-t_0)},
\]

so fixed-budget AI loss strictly decreases over time. This proves a precise version of claim $(i)$.

However, claim $(ii)$ is not well-defined without choosing a capability index, and different legitimate choices give decreasing, constant, or increasing rates of gain.

Claim $(iii)$ is not derivable from $A1$ because economic dominance requires additional assumptions about payoffs, costs, preferences, and adoption constraints.

Therefore, the stated theorem $C1$ cannot be rigorously proved using only the quantitative content of $A1$.